\documentclass[../../../uem_tok_q.tex]{subfiles}

\begin{document}
    複素数平面上の点$\myfraca{1}{2}$を中心とする半径$\myfraca{1}{2}$%
    の円の周から原点を除いた曲線を$C$とする．%

    \begin{enumerate}
        \item 曲線$C$上の複素数$z$に対し，$\myfraca{1}{z}$の実部は1であることを示せ．
        \item $\alpha$，$\beta$を曲線$C$上の相異なる複素数とするとき，%
            $\myfraca{1}{\alpha^2}+\myfraca{1}{\beta^2}$がとりうる範囲を複素数平面上に図示せよ．%
            \label{uem_tok_2025_f_sci_06_q:2}
        \item $\gamma$を\,\ref{uem_tok_2025_f_sci_06_q:2}\,で求めた範囲に属さない複素数とするとき，%
            $\myfraca{1}{\gamma}$の実部がとりうる値の最大値と最小値を求めよ．
    \end{enumerate}
\end{document}
